\chapter{Tabellen zum Ausfüllen}

Füllen Sie diese Seiten von Hand aus, solange Sie das System vor sich haben. Später ist das die einzige Stelle, an der steht, was Sie eingerichtet haben.
\renewcommand{\arraystretch}{1.9}

\section{Hardware}
\begin{tabularx}{\linewidth}{@{}>{\RaggedRight\arraybackslash}p{4.8cm}X@{}}
\toprule
Rechner (Hersteller, Modell) & \\
\midrule
Arbeitsspeicher (GB) & \\
\midrule
USB-SSD (Hersteller, Größe) & \\
\midrule
Startmenü-Taste & \\
\midrule
Secure Boot im BIOS geändert? (ja/nein) & \\
\midrule
Netzteil / Ladegerät (Watt) & \\
\midrule
Wo liegt das Zubehör? & \\
\bottomrule
\end{tabularx}

\section{Softwarestand}
\begin{tabularx}{\linewidth}{@{}>{\RaggedRight\arraybackslash}p{4.4cm}XX@{}}
\toprule
 & \textbf{Beim Einrichten} & \textbf{Letzte Prüfung} \\
\midrule
Datum & & \\
\midrule
NOMAD-Version (Fußzeile der Startseite) & & \\
\midrule
Ubuntu-Version & & \\
\midrule
Image-Datei / Prüfsumme & & \\
\midrule
Letztes Update eingespielt am & & \\
\bottomrule
\end{tabularx}

\hinweis{Wie ein Update geht: Startseite $\rightarrow$ \menue{Einstellungen} $\rightarrow$ \menue{System-Update} $\rightarrow$ „Erneut prüfen“, bei Bedarf „Update starten“. Dafür braucht der Rechner Internet. Danach der Offline-Test (Kapitel~6).}

\section{Zugangsdaten}
\achtung{Dieses Heft ist ein Papier. Schreiben Sie nur Passwörter hinein, wenn es an einem sicheren Ort liegt.}
\begin{tabularx}{\linewidth}{@{}>{\RaggedRight\arraybackslash}p{4.8cm}X@{}}
\toprule
Benutzername des Systems & \\
\midrule
Passwort des Systems & \\
\midrule
WLAN-Name & \\
\midrule
WLAN-Passwort & \\
\midrule
Router-Adresse / Passwort & \\
\midrule
Adresse von NOMAD im Heimnetz & \\
\bottomrule
\end{tabularx}

\section{Geladene Inhalte}
\begin{tabularx}{\linewidth}{@{}>{\RaggedRight\arraybackslash}p{4.4cm}>{\RaggedRight\arraybackslash}p{2.2cm}X@{}}
\toprule
\textbf{Inhalt} & \textbf{Größe} & \textbf{Geladen am / Anmerkung} \\
\midrule
Wikipedia (Auswahl) & & \\
\midrule
Medizin & & \\
\midrule
Überleben und Vorsorge & & \\
\midrule
Bildung und Nachschlagen & & \\
\midrule
Heimwerken und Reparatur & & \\
\midrule
Landwirtschaft und Essen & & \\
\midrule
Computer und Technik & & \\
\midrule
Karten & & \\
\midrule
KI-Modell & & \\
\bottomrule
\end{tabularx}

\section{Sicherung und Offline-Test}
\begin{tabularx}{\linewidth}{@{}>{\RaggedRight\arraybackslash}p{3.2cm}XXX@{}}
\toprule
\textbf{Datum} & \textbf{Offline-Test bestanden?} & \textbf{Sicherung (Wo?)} & \textbf{Unterschrift} \\
\midrule
 & & & \\
\midrule
 & & & \\
\midrule
 & & & \\
\midrule
 & & & \\
\bottomrule
\end{tabularx}

\hinweis{Eine Sicherung (Backup) schützt gegen einen Ausfall der USB-SSD. Am einfachsten: Eine zweite USB-SSD mit demselben System (Kapitel~2) herstellen und die Inhalte nachladen. Das Heft hilft Ihnen dabei, den Stand nachzuvollziehen.}
